% GENERATED FILE. Edit canonical lessons, then run npm run generate:books.
% canonical-source-hash: fnv1a64:044580ed35336269
% canonical-lessons: FR-C203-defaite, FR-C203-championnat, FR-C203-joueur, FR-C203-entraineur, FR-C203-correspondance

\chapter{Defeat, Championship, Player, Coach, Connection}
\label{ch:fr-defeat-championship-player-coach-connection}
\hlchaptermodality{203}

\begin{chapteropening}
\textbf{By the end of this chapter:} \emph{I can say a defeat, a championship, a player, a coach and a connection.}

Complete the last lesson of chapter 203: I can say a defeat, a championship, a player, a coach and a connection.
\end{chapteropening}

\section[la défaite]{la défaite — a defeat}

\label{lesson:FR-C203-defaite}



\begin{quote}
Before the new one: say the French for a small bag, then the French for a rope.
\end{quote}

\subsection*{You'll want to know: la défaite}
\textbf{la défaite} — \textquotedblleft{}a defeat\textquotedblright{}.

\begin{cousinweb}
\textbf{défaite} is the past participle of the verb for \textquotedblleft{}to undo\textquotedblright{}, used as a noun: an undoing. Old French had \emph{desfait}, \textquotedblleft{}undone\textquotedblright{} (EtymDB 2.1); the verb is built on Latin \emph{facere}, \textquotedblleft{}to do\textquotedblright{}.
\end{cousinweb}

\begin{grammarlens}[title={What you've built}]
One more word for this chapter.
\end{grammarlens}

\subsection*{Your turn}
\begin{itemize}
\raggedright
  \item \emph{Say it:} \emph{la défaite}
  \item \emph{Say it:} \emph{la défaite}, once more
  \item \emph{Say it:} say \emph{la défaite}
  \item \emph{Recall:} say the French for a small bag, then the French for a rope, then say \emph{la défaite} again
\end{itemize}

\begin{tcolorbox}[breakable,colback=teal!4,colframe=teal!35!black,title={Before you move on}]
What does la défaite mean? (\textquotedblleft{}A defeat\textquotedblright{}.) Say it once more.
\end{tcolorbox}
\section[le championnat]{le championnat — a championship}

\label{lesson:FR-C203-championnat}



\begin{quote}
Before the new one: say the French for a rope, then the French for a defeat.
\end{quote}

\subsection*{You'll want to know: le championnat}
\textbf{le championnat} — \textquotedblleft{}a championship\textquotedblright{}.

\begin{cousinweb}
\textbf{championnat} — built from the word for \textquotedblleft{}champion\textquotedblright{}, $\leftarrow$ Medieval Latin \emph{campio}, \textquotedblleft{}a fighter in a duel\textquotedblright{}, which EtymDB 2.1 takes from Frankish \emph{*kampjo}, \textquotedblleft{}a fighter\textquotedblright{}, with \emph{-at}.
\end{cousinweb}

\begin{grammarlens}[title={What you've built}]
One more word for this chapter.
\end{grammarlens}

\subsection*{Your turn}
\begin{itemize}
\raggedright
  \item \emph{Say it:} \emph{le championnat}
  \item \emph{Say it:} \emph{le championnat}, once more
  \item \emph{Say it:} say \emph{le championnat}
  \item \emph{Recall:} say the French for a rope, then the French for a defeat, then say \emph{le championnat} again
\end{itemize}

\begin{tcolorbox}[breakable,colback=teal!4,colframe=teal!35!black,title={Before you move on}]
What does le championnat mean? (\textquotedblleft{}A championship\textquotedblright{}.) Say it once more.
\end{tcolorbox}
\section[le joueur]{le joueur — a player}

\label{lesson:FR-C203-joueur}



\begin{quote}
Before the new one: say the French for a defeat, then the French for a championship.
\end{quote}

\subsection*{You'll want to know: le joueur}
\textbf{le joueur} — \textquotedblleft{}a player\textquotedblright{}.

\begin{cousinweb}
\textbf{joueur} — \emph{jouer}, \textquotedblleft{}to play\textquotedblright{}, with \emph{-eur}, \textquotedblleft{}one who\textquotedblright{}.
\end{cousinweb}

\begin{grammarlens}[title={What you've built}]
One more word for this chapter.
\end{grammarlens}

\subsection*{Your turn}
\begin{itemize}
\raggedright
  \item \emph{Say it:} \emph{le joueur}
  \item \emph{Say it:} \emph{le joueur}, once more
  \item \emph{Say it:} say \emph{le joueur}
  \item \emph{Recall:} say the French for a defeat, then the French for a championship, then say \emph{le joueur} again
\end{itemize}

\begin{tcolorbox}[breakable,colback=teal!4,colframe=teal!35!black,title={Before you move on}]
What does le joueur mean? (\textquotedblleft{}A player\textquotedblright{}.) Say it once more.
\end{tcolorbox}
\section[l'entraîneur]{l'entraîneur — a coach}

\label{lesson:FR-C203-entraineur}



\begin{quote}
Before the new one: say the French for a championship, then the French for a player.
\end{quote}

\subsection*{You'll want to know: l'entraîneur}
\textbf{l'entraîneur} — \textquotedblleft{}a coach\textquotedblright{}.

\begin{cousinweb}
\textbf{entraîneur} — built from the verb for \textquotedblleft{}to train, to draw along\textquotedblright{}, from \emph{en-} and the verb for \textquotedblleft{}to drag\textquotedblright{}, $\leftarrow$ Vulgar Latin \emph{*traginare}, from Latin \emph{trahere}, \textquotedblleft{}to draw\textquotedblright{} (Lewis \& Short).
\end{cousinweb}

\begin{grammarlens}[title={What you've built}]
One more word for this chapter.
\end{grammarlens}

\subsection*{Your turn}
\begin{itemize}
\raggedright
  \item \emph{Say it:} \emph{l'entraîneur}
  \item \emph{Say it:} \emph{l'entraîneur}, once more
  \item \emph{Say it:} say \emph{l'entraîneur}
  \item \emph{Recall:} say the French for a championship, then the French for a player, then say \emph{l'entraîneur} again
\end{itemize}

\begin{tcolorbox}[breakable,colback=teal!4,colframe=teal!35!black,title={Before you move on}]
What does l'entraîneur mean? (\textquotedblleft{}A coach\textquotedblright{}.) Say it once more.
\end{tcolorbox}
\section[la correspondance]{la correspondance — a connection}

\label{lesson:FR-C203-correspondance}



\begin{quote}
Before the new one: say the French for a player, then the French for a coach.
\end{quote}

\subsection*{You'll want to know: la correspondance}
\textbf{la correspondance} — \textquotedblleft{}a connection\textquotedblright{}.

\begin{cousinweb}
\textbf{correspondance} — built from the verb for \textquotedblleft{}to correspond\textquotedblright{}, $\leftarrow$ Medieval Latin \emph{correspondere}, from \emph{com-} and \emph{respondere}, \textquotedblleft{}to answer\textquotedblright{}.
\end{cousinweb}

\begin{grammarlens}[title={What you've built}]
One more word for this chapter.
\end{grammarlens}

\subsection*{Your turn}
\begin{itemize}
\raggedright
  \item \emph{Say it:} \emph{la correspondance}
  \item \emph{Say it:} \emph{la correspondance}, once more
  \item \emph{Say it:} say \emph{la correspondance}
  \item \emph{Recall:} say the French for a player, then the French for a coach, then say \emph{la correspondance} again
\end{itemize}

\begin{tcolorbox}[breakable,colback=teal!4,colframe=teal!35!black,title={Before you move on}]
What does la correspondance mean? (\textquotedblleft{}A connection\textquotedblright{}.) Say it once more.
\end{tcolorbox}
